\section{Conclusion}

A two-dimensional finite-volume solver for the steady-state heat conduction equation was developed and verified. The Laplace equation was discretized using a second-order cell-centered finite-volume formulation on a structured Cartesian mesh. Four solution methods were implemented: Point Gauss-Seidel, Successive Over-Relaxation (SOR), Line-SOR, and a fully implicit direct method. Dirichlet boundary conditions were imposed through ghost cells, allowing a consistent treatment of the physical boundaries throughout the computational domain.

The numerical solutions were assessed through comparison with the analytical solution using temperature contours, line-cut comparisons, and error distributions. The fully implicit solution on a $40\times40$ mesh exhibited an $L_2$ error of $1.6122\times10^{-4}$ and a maximum absolute error of $7.2049\times10^{-4}$. Residual histories and final residual values were also used to assess the convergence of the numerical methods and their consistency with the discretized governing equation.

The performance studies demonstrated the influence of the solution algorithm and relaxation factor on convergence and computational cost. For the tested $40\times40$ benchmark with $\omega=1.1$, Line-SOR required approximately half as many iterations as point-wise SOR and achieved a substantially lower recorded CPU time. The SOR comparisons also showed that increasing the relaxation factor from $1.1$ to $1.5$ reduced the iteration count, although the computational benefit depended on the selected factor. Across the tested mesh sizes from $N=10$ to $N=40$, the fully implicit direct method had the lowest recorded CPU time. However, its dense coefficient matrix imposes increasing memory and computational costs as the problem size grows, limiting the extent to which this result can be generalized to larger meshes.

The modular software architecture separates numerical solvers, shared computational functions, study definitions, and visualization. This organization facilitates systematic comparisons between numerical methods and provides a foundation for further development. Future work may include more extensive mesh-convergence and performance studies, more efficient solution of the tridiagonal systems arising in Line-SOR, and the implementation of additional numerical methods, such as sparse direct solvers and multigrid techniques. The framework can also be extended to more general transport equations encountered in computational fluid dynamics.

Overall, this project establishes a verified computational framework for studying finite-volume discretization and numerical solution strategies. It also provides a basis for further investigating the relationship between numerical accuracy, convergence behavior, and computational efficiency.